% Template da Introdução — regras diferentes do template genérico
% de seção (assets/template-section.tex). Ver detalhes em
% references/03-estrutura-taxonomias.md, seção "Regras de conteúdo
% da Introdução". Resumo das restrições, para não esquecer ao
% preencher:
%   - SEM \subsection{} — a introdução é um bloco de prosa corrido,
%     nunca dividida em subseções.
%   - SEM figuras, tabelas, gráficos ou \includegraphics/\begin{table}
%     de qualquer tipo — a primeira ilustração do artigo só aparece
%     a partir da seção 2 (Background, se houver, ou a subseção de
%     taxonomia do corpo principal).
%   - NÃO detalhar o framework/taxonomia adotada — apenas resumir a
%     visão geral em 1-2 frases. O detalhamento completo é exclusivo
%     da primeira subseção do corpo principal (taxonomia adotada).
%   - NÃO aprofundar o estado da arte — isso é papel da seção de
%     Background (quando houver) ou dos eixos do corpo principal,
%     não da introdução.
%   - Trabalhos relacionados/surveys existentes entram de forma
%     SUPERFICIAL: o suficiente para posicionar a contribuição desta
%     survey (o que ela oferece que as outras não oferecem), sem
%     comparação detalhada, tabela ou crítica aprofundada — essa
%     comparação detalhada (inclusive a tabela comparativa com
%     surveys relacionadas, se houver) vai para o Background ou para
%     a subseção de taxonomia do corpo principal, nunca aqui.
\section{Introdução}
% TODO (Fase 3/4): escrever em prosa corrida, cobrindo nesta ordem:
%   1. Motivação e contexto do tema (por que ele importa agora).
%   2. Escopo da survey (o que está dentro/fora, reaproveitando a
%      crítica de tema da Fase 1).
%   3. Posicionamento superficial frente a surveys existentes —
%      1-2 frases citando as principais surveys relacionadas
%      (\cite{}) e o que esta survey acrescenta, SEM tabela e SEM
%      comparação detalhada.
%   4. Resumo de 1-2 frases da visão geral do framework/taxonomia
%      adotada (não o detalhamento — isso é da primeira subseção do
%      corpo principal).
%   5. Contribuições explícitas em lista (\begin{itemize}), objetivas.
%   6. Roteiro do restante do artigo (1 frase por seção).
